\documentclass[10pt,a4paper]{article}

%double si on veut une version prof et une version élève. simple si on veut une seule version.
\newcommand{\buildmode}{simple}

\InputIfFileExists{version.tex}{}{}
% Version (définie par le script)
\providecommand{\version}{eleve}

\usepackage{feuilleinterro}

%%%à modifier à chaque fois

\renewcommand{\niveau}{1G}
\renewcommand{\chapitre}{Probabilités conditionnelles}
\renewcommand{\numchap}{2}
\renewcommand{\theme}{Probabilités}

%%%titre "CORRECTION" au lieu de "INTERROGATION", sans les champs Nom/Prénom
\renewcommand{\interroversion}[1]{
    \noindent\textbf{\niveau}

    \vspace{-0.5cm}
    \begin{center}
        \textbf{\large CORRECTION -- \chapitre}
    \end{center}

    \vspace{0.1cm}
}

%%%arbre pondéré à deux niveaux
%%% \arbreproba{A}{B}{P(A)}{P(Abar)}{P_A(B)}{P_A(Bbar)}{P_Abar(B)}{P_Abar(Bbar)}
\newcommand{\arbreproba}[8]{%
\begin{tikzpicture}[
  grow'=right,
  level distance=3.6cm,
  level 1/.style={sibling distance=2.4cm},
  level 2/.style={sibling distance=1.2cm},
  edge from parent/.style={draw,-latex},
  every node/.style={font=\small}
]
\node {}
  child { node {\(#1\)}
    child { node {\(#2\)} edge from parent node[above] {#5} }
    child { node {\(\overline{#2}\)} edge from parent node[below] {#6} }
    edge from parent node[above] {#3}
  }
  child { node {\(\overline{#1}\)}
    child { node {\(#2\)} edge from parent node[above] {#7} }
    child { node {\(\overline{#2}\)} edge from parent node[below] {#8} }
    edge from parent node[below] {#4}
  };
\end{tikzpicture}%
}
\newcommand{\bl}[1]{\textcolor{blue}{\(#1\)}}

\renewcommand{\arraystretch}{1.4}
\newcolumntype{C}{>{\centering\arraybackslash}p{2.4cm}}

\begin{document}

% =========================================================
% PAGE 1 : VERSION A
% =========================================================

\interroversion{A}

\textbf{Exercice 1 -- Calculs avec des fractions \hfill /2}

Pour \(x\neq2\), on met les deux fractions au même dénominateur :
\[
A=\dfrac{2}{3}+\dfrac{x+1}{x-2}=\dfrac{2(x-2)}{3(x-2)}+\dfrac{3(x+1)}{3(x-2)}=\dfrac{2x-4+3x+3}{3(x-2)}=\textcolor{blue}{\dfrac{5x-1}{3(x-2)}}.
\]

\vspace{0.3cm}

\textbf{Exercice 2 -- Demi-pension et sport \hfill /8}

\begin{enumerate}
\item Calculs utilisant les pourcentages :
\begin{itemize}
\item demi-pensionnaires : \(\dfrac{60}{100}\times200=120\), donc externes : \(200-120=80\) ;
\item sport en club : \(\dfrac{31}{100}\times200=62\), donc pas de sport : \(200-62=138\).
\end{itemize}
On complète ensuite par différences : \(120-90=30\) et \(80-32=48\) (vérification : \(90+48=138\)).

\begin{center}
\begin{tabular}{|l|C|C|C|}
    \hline
     & Sport en club & Pas de sport & Total \\
    \hline
    Demi-pensionnaires & \bl{30} & 90 & \bl{120} \\
    \hline
    Externes & 32 & \bl{48} & \bl{80} \\
    \hline
    Total & \bl{62} & \bl{138} & 200 \\
    \hline
\end{tabular}
\end{center}

\item \(P(D)=\dfrac{120}{200}=\textcolor{blue}{0{,}6}\).

\item \(D\cap S\) : \og l'élève est demi-pensionnaire \textbf{et} pratique un sport en club \fg{}. \(P(D\cap S)=\dfrac{30}{200}=\textcolor{blue}{0{,}15}\).

\item Parmi les \(120\) demi-pensionnaires, \(30\) pratiquent un sport en club : \(P_D(S)=\dfrac{30}{120}=\textcolor{blue}{0{,}25}\).

\item \(P(\overline D)=\dfrac{80}{200}=0{,}4\) ; \(P_D(\overline S)=\dfrac{90}{120}=0{,}75\) ; \(P_{\overline D}(S)=\dfrac{32}{80}=0{,}4\) ; \(P_{\overline D}(\overline S)=\dfrac{48}{80}=0{,}6\).

\begin{center}
\arbreproba{D}{S}{\bl{0{,}6}}{\bl{0{,}4}}{\bl{0{,}25}}{\bl{0{,}75}}{\bl{0{,}4}}{\bl{0{,}6}}
\end{center}

\item On multiplie les probabilités le long du chemin \(D\to S\) :
\[
P(D\cap S)=P(D)\times P_D(S)=0{,}6\times0{,}25=\textcolor{blue}{0{,}15}.
\]
On retrouve bien le résultat de la question 3.
\end{enumerate}

\newpage

% =========================================================
% PAGE 2 : VERSION B
% =========================================================

\interroversion{B}

\textbf{Exercice 1 -- Calculs avec des fractions \hfill /2}

Pour \(x\neq-3\), on met les deux fractions au même dénominateur :
\[
A=\dfrac{3}{4}+\dfrac{x-1}{x+3}=\dfrac{3(x+3)}{4(x+3)}+\dfrac{4(x-1)}{4(x+3)}=\dfrac{3x+9+4x-4}{4(x+3)}=\textcolor{blue}{\dfrac{7x+5}{4(x+3)}}.
\]

\vspace{0.3cm}

\textbf{Exercice 2 -- Au cinéma \hfill /8}

\begin{enumerate}
\item Calculs utilisant les pourcentages :
\begin{itemize}
\item tarif réduit : \(\dfrac{40}{100}\times400=160\), donc plein tarif : \(400-160=240\) ;
\item film en 3D : \(\dfrac{27}{100}\times400=108\), donc film en 2D : \(400-108=292\).
\end{itemize}
On complète ensuite par différences : \(160-112=48\) et \(240-60=180\) (vérification : \(112+180=292\)).

\begin{center}
\begin{tabular}{|l|C|C|C|}
    \hline
     & Film en 3D & Film en 2D & Total \\
    \hline
    Tarif réduit & \bl{48} & 112 & \bl{160} \\
    \hline
    Plein tarif & 60 & \bl{180} & \bl{240} \\
    \hline
    Total & \bl{108} & \bl{292} & 400 \\
    \hline
\end{tabular}
\end{center}

\item \(P(R)=\dfrac{160}{400}=\textcolor{blue}{0{,}4}\).

\item \(R\cap T\) : \og le spectateur a payé un tarif réduit \textbf{et} a vu un film en 3D \fg{}. \(P(R\cap T)=\dfrac{48}{400}=\textcolor{blue}{0{,}12}\).

\item Parmi les \(160\) spectateurs au tarif réduit, \(48\) ont vu un film en 3D : \(P_R(T)=\dfrac{48}{160}=\textcolor{blue}{0{,}3}\).

\item \(P(\overline R)=\dfrac{240}{400}=0{,}6\) ; \(P_R(\overline T)=\dfrac{112}{160}=0{,}7\) ; \(P_{\overline R}(T)=\dfrac{60}{240}=0{,}25\) ; \(P_{\overline R}(\overline T)=\dfrac{180}{240}=0{,}75\).

\begin{center}
\arbreproba{R}{T}{\bl{0{,}4}}{\bl{0{,}6}}{\bl{0{,}3}}{\bl{0{,}7}}{\bl{0{,}25}}{\bl{0{,}75}}
\end{center}

\item On multiplie les probabilités le long du chemin \(R\to T\) :
\[
P(R\cap T)=P(R)\times P_R(T)=0{,}4\times0{,}3=\textcolor{blue}{0{,}12}.
\]
On retrouve bien le résultat de la question 3.
\end{enumerate}

\newpage

% =========================================================
% PAGE 3 : VERSION C
% =========================================================

\interroversion{C}

\textbf{Exercice 1 -- Calculs avec des fractions \hfill /2}

Pour \(x\neq1\), on met les deux fractions au même dénominateur :
\[
A=\dfrac{1}{5}+\dfrac{2x+1}{x-1}=\dfrac{x-1}{5(x-1)}+\dfrac{5(2x+1)}{5(x-1)}=\dfrac{x-1+10x+5}{5(x-1)}=\textcolor{blue}{\dfrac{11x+4}{5(x-1)}}.
\]

\vspace{0.3cm}

\textbf{Exercice 2 -- Clients d'un magasin \hfill /8}

\begin{enumerate}
\item Calculs utilisant les pourcentages :
\begin{itemize}
\item avec carte : \(\dfrac{60}{100}\times500=300\), donc sans carte : \(500-300=200\) ;
\item plus de \(50\) \euro{} : \(\dfrac{30}{100}\times500=150\), donc \(50\) \euro{} ou moins : \(500-150=350\).
\end{itemize}
On complète ensuite par différences : \(300-180=120\) et \(200-30=170\) (vérification : \(180+170=350\)).

\begin{center}
\begin{tabular}{|l|C|C|C|}
    \hline
     & Plus de 50 \euro & 50 \euro{} ou moins & Total \\
    \hline
    Avec carte & \bl{120} & 180 & \bl{300} \\
    \hline
    Sans carte & 30 & \bl{170} & \bl{200} \\
    \hline
    Total & \bl{150} & \bl{350} & 500 \\
    \hline
\end{tabular}
\end{center}

\item \(P(F)=\dfrac{300}{500}=\textcolor{blue}{0{,}6}\).

\item \(F\cap M\) : \og le client possède la carte \textbf{et} a fait un achat de plus de \(50\) \euro{} \fg{}. \(P(F\cap M)=\dfrac{120}{500}=\textcolor{blue}{0{,}24}\).

\item Parmi les \(300\) clients ayant la carte, \(120\) ont fait un achat de plus de \(50\) \euro{} : \(P_F(M)=\dfrac{120}{300}=\textcolor{blue}{0{,}4}\).

\item \(P(\overline F)=\dfrac{200}{500}=0{,}4\) ; \(P_F(\overline M)=\dfrac{180}{300}=0{,}6\) ; \(P_{\overline F}(M)=\dfrac{30}{200}=0{,}15\) ; \(P_{\overline F}(\overline M)=\dfrac{170}{200}=0{,}85\).

\begin{center}
\arbreproba{F}{M}{\bl{0{,}6}}{\bl{0{,}4}}{\bl{0{,}4}}{\bl{0{,}6}}{\bl{0{,}15}}{\bl{0{,}85}}
\end{center}

\item On multiplie les probabilités le long du chemin \(F\to M\) :
\[
P(F\cap M)=P(F)\times P_F(M)=0{,}6\times0{,}4=\textcolor{blue}{0{,}24}.
\]
On retrouve bien le résultat de la question 3.
\end{enumerate}

\newpage

% =========================================================
% PAGE 4 : VERSION D
% =========================================================

\interroversion{D}

\textbf{Exercice 1 -- Calculs avec des fractions \hfill /2}

Pour \(x\neq-1\), on met les deux fractions au même dénominateur :
\[
A=\dfrac{3}{2}+\dfrac{x+2}{x+1}=\dfrac{3(x+1)}{2(x+1)}+\dfrac{2(x+2)}{2(x+1)}=\dfrac{3x+3+2x+4}{2(x+1)}=\textcolor{blue}{\dfrac{5x+7}{2(x+1)}}.
\]

\vspace{0.3cm}

\textbf{Exercice 2 -- Une course à pied \hfill /8}

\begin{enumerate}
\item Calculs utilisant les pourcentages :
\begin{itemize}
\item licenciés : \(\dfrac{40}{100}\times250=100\), donc non licenciés : \(250-100=150\) ;
\item moins d'une heure : \(\dfrac{36}{100}\times250=90\), donc une heure ou plus : \(250-90=160\).
\end{itemize}
On complète ensuite par différences : \(100-40=60\) et \(150-30=120\) (vérification : \(40+120=160\)).

\begin{center}
\begin{tabular}{|l|C|C|C|}
    \hline
     & Moins d'1 h & 1 h ou plus & Total \\
    \hline
    Licenciés & \bl{60} & 40 & \bl{100} \\
    \hline
    Non licenciés & 30 & \bl{120} & \bl{150} \\
    \hline
    Total & \bl{90} & \bl{160} & 250 \\
    \hline
\end{tabular}
\end{center}

\item \(P(L)=\dfrac{100}{250}=\textcolor{blue}{0{,}4}\).

\item \(L\cap H\) : \og le coureur est licencié \textbf{et} a terminé en moins d'une heure \fg{}. \(P(L\cap H)=\dfrac{60}{250}=\textcolor{blue}{0{,}24}\).

\item Parmi les \(100\) licenciés, \(60\) ont terminé en moins d'une heure : \(P_L(H)=\dfrac{60}{100}=\textcolor{blue}{0{,}6}\).

\item \(P(\overline L)=\dfrac{150}{250}=0{,}6\) ; \(P_L(\overline H)=\dfrac{40}{100}=0{,}4\) ; \(P_{\overline L}(H)=\dfrac{30}{150}=0{,}2\) ; \(P_{\overline L}(\overline H)=\dfrac{120}{150}=0{,}8\).

\begin{center}
\arbreproba{L}{H}{\bl{0{,}4}}{\bl{0{,}6}}{\bl{0{,}6}}{\bl{0{,}4}}{\bl{0{,}2}}{\bl{0{,}8}}
\end{center}

\item On multiplie les probabilités le long du chemin \(L\to H\) :
\[
P(L\cap H)=P(L)\times P_L(H)=0{,}4\times0{,}6=\textcolor{blue}{0{,}24}.
\]
On retrouve bien le résultat de la question 3.
\end{enumerate}

\end{document}
